\title{Parameter-Efficient Orthogonal Finetuning via Factorization}

\begin{abstract}
The widespread adoption of large foundation models highlights the prohibitive computational costs associated with training such systems from the ground up. Consequently, there is a growing need for methodologies that can adapt these pre-trained models to specific downstream applications in an efficient manner. In this work, we revisit a structured approach for model adaptation, known as Orthogonal Finetuning (OFT), which enables effective downstream performance. However, OFT typically involves a substantial number of tunable parameters due to the inherently high dimensionality of orthogonal transformations. To tackle this, we first reinterpret OFT through the lens of information transmission, identifying critical criteria for enhancing parameter efficiency. Drawing inspiration from the Cooley-Tukey fast Fourier transform algorithm’s capability for rapid information transfer, we introduce an efficient orthogonal parameterization method utilizing butterfly networks. We embed this structure into OFT, presenting a new parameter-efficient finetuning technique called Orthogonal Butterfly~(BOFT). BOFT generalizes the OFT paradigm, with the latter as a particular instance within its broader framework. Comprehensive experimental evaluations are conducted, demonstrating BOFT’s efficacy in customizing large vision transformers, expansive language models, and text-to-image diffusion systems for a variety of downstream vision and language tasks.
\end{abstract}

\section{Introduction}

Recent advancements exemplified by models like ChatGPT~\cite{brown2020language,chen2021evaluating} and Stable Diffusion~\cite{rombach2022high} have showcased the exceptional transfer capabilities of large foundation models. This impressive increase in capability has coincided with a dramatic escalation in the models’ parameter counts (\emph{e.g.}, GPT-3 features approximately 175 billion parameters), making the task of training such architectures from scratch increasingly infeasible. To foster ongoing progress in the field, it has become essential to develop strategies that allow for adapting pre-existing foundation models rather than undertaking resource-intensive retraining. Efficient adaptation approaches are thus vital for leveraging the power of existing models on new tasks. The prevailing strategies for effective model adaptation include: \emph{model finetuning}~\cite{hulora2022,zaken2022bitfit,guo2020parameter,raffel2020exploring,qiu2023controlling,zhang2022adaptive,chavan2023one}, which selectively updates a portion of the model’s parameters while maintaining its original structure; \emph{adapter tuning}~\cite{rebuffi2017learning,houlsby2019parameter,pfeiffer2020adapterfusion,lin2020exploring,he2021towards}, which incorporates additional trainable modules into the network; and \emph{prompt tuning}~\cite{li2021prefix,lester2021power}, which prepends trainable prompts to the model’s input. Of these techniques, model finetuning stands out for its simplicity, strong empirical performance, and the advantage of not increasing inference latency.

The central idea underlying model finetuning is to maintain the similarity between the pretrained and finetuned models according to specific metrics, thereby retaining the pretraining information. In practice, many prevalent finetuning techniques employ a low learning rate as an informal means to limit the divergence between the pretrained and updated models. Recognizing that weight matrices are inherently structured, a more systematic strategy involves preserving their relational characteristics, such as the pairwise angular relationships between neurons. This consideration forms the basis of Orthogonal Finetuning~(OFT)~\cite{qiu2023controlling}, a recently developed framework in which all neurons within a given layer undergo transformation via a common orthogonal matrix, thereby rigorously maintaining pairwise angular distances throughout finetuning. While OFT has shown advantageous properties with regard to generalization and convergence in finetuning tasks for text-to-image diffusion models~\cite{qiu2023controlling}, the approach can lead to a substantial number of trainable parameters, primarily because orthogonal matrices in high dimensions are large. To enhance parameter efficiency, OFT employs a block-diagonal design, significantly lowering parameter counts. Nonetheless, this increase in efficiency introduces drawbacks—the orthogonal matrices possess a rigid sparsity pattern and each block operates in isolation, with independent orthogonal transformations. Such block-diagonal sparsity, while economical in terms of parameters, may lead to undesirable biases (for instance, the reduced expressiveness of block-diagonal orthogonal matrices precludes them from representing many conventional linear transformations), and crucially, the task of discovering optimal sparsity patterns remains unsolved.

Addressing this issue requires constructing a dense orthogonal matrix in a manner that remains efficient with respect to the number of parameters. At first glance, this appears challenging because a $d$-dimensional dense orthogonal matrix typically involves $\mathcal{O}(d^2)$ parameters. However, our work introduces a new methodology, wherein a dense orthogonal matrix is built through the composition of multiple sparse matrices in a factorized fashion. The underlying principle is that parameter savings can be realized by opting for increased computational cost in exchange for a reduced memory footprint. Since our model expresses the orthogonal matrix as a product of sparse factors, this multiplication must be repeatedly calculated at every iteration during training. To contextualize our matrix factorization, we reframe the task of producing a dense orthogonal matrix as one of transmitting information over a network. With this perspective, constructing a dense orthogonal matrix via products of sparse matrices can be conceptualized as communicating information through a graph with grid-like topology. This lens enables us to design diverse approaches to sparse matrix factorization that retain sufficient representational power for dense matrices while tightly controlling the number of tunable parameters.

To enhance parameter efficiency within this information transmission setting, we take inspiration from the butterfly network architecture found in the Cooley-Tukey fast Fourier transform algorithm~\cite{cooley1965algorithm}, a structure that facilitates the rapid exchange of information among nodes~\cite{klasing1994broadcasting}. The butterfly pattern inherent in the Cooley-Tukey algorithm offers a natural blueprint for a type of sparse matrix factorization, known as butterfly factorization. Leveraging this, it becomes possible to express a $d$-dimensional dense matrix as a product of $\mathcal{O}(\log d)$ sparse matrices, where each matrix contains only $\mathcal{O}(d)$ nonzero elements. By guaranteeing that each component matrix remains orthogonal, our strategy results in an orthogonal parameterization that requires just $\mathcal{O}(d\log d)$ parameters—dramatically fewer than a fully dense orthogonal matrix. Building on this efficient parametrization, we introduce Orthogonal Butterfly~(BOFT), a new parameter-efficient finetuning technique. BOFT generalizes OFT, offering a comprehensive orthogonal finetuning framework. Notably, both block-diagonal and butterfly constructions share a key property: they introduce sparsity to orthogonal matrices, thereby constraining the number of tunable parameters. A comparison with the low-rank decomposition used in LoRA~\cite{hulora2022} highlights that both sparse and low-rank matrices belong to principal categories of structured matrices~\cite{candes2011robust} used to attain parameter efficiency.

In contrast to the block-diagonal form utilized in OFT to balance between expressiveness and structural regularity, BOFT employs a butterfly configuration, which allows for a more continuous transition between matrices from the full orthogonal group (capturing expressiveness) and the identity matrix (favoring regularity). This facilitates the identification of a more compact hypothesis space within the orthogonal group for downstream applications. Given that butterfly architectures are central to a range of efficient linear transforms, such as the discrete Fourier transform and the discrete cosine transform, we posit that this structure introduces an inherent inductive bias which promotes generalization and guards against overfitting. Our perspective is shaped by the observation that, whereas the butterfly structure readily accommodates many well-known linear transforms, the block-diagonal construction employed in OFT lacks this ability. Our principal contributions are summarized as follows:

\begin{itemize}[leftmargin=*,nosep]
    \item We examine the parameter efficiency challenge in orthogonal finetuning from the perspective of an information transmission framework, recasting the construction of a parameter-efficient dense orthogonal matrix as an information flow problem in a grid-like graph structure.
    \item Drawing on butterfly patterns akin to those in the Cooley-Tukey algorithm, we introduce Orthogonal Butterfly, an orthogonal finetuning technique optimized for parameter efficiency, developed within this information transmission paradigm.
    \item We offer several theoretical explanations for how BOFT achieves substantial reduction in trainable parameters without compromising expressive capacity or training robustness. Owing to its matrix factorization, BOFT also admits an appealing mechanism for weight interpolation (refer to Figure~\ref{fig:boft_interpolation}).
    \item We are the first to deploy orthogonal finetuning for a broad array of tasks beyond controllable text-to-image generation~\cite{qiu2023controlling}, thereby highlighting its versatility as a general-purpose finetuning strategy. In this effort, we apply BOFT across diverse adaptation tasks in domains spanning computer vision and natural language processing. BOFT consistently surpasses existing state-of-the-art approaches by a wide margin, demonstrating its advantages in parameter efficiency and generalizability.
\end{itemize}

\section{Related Work}

\textbf{Parameter-efficient finetuning (PEFT)}. The emergence of increasingly large and capable foundation models (\emph{e.g.}, \cite{brown2020language,radford2021learning,rombach2022high,kirillov2023segment}) has fueled substantial research into parameter-efficient approaches for finetuning such models on downstream applications~\cite{treviso2023efficient,ding2023parameter,lialin2023scaling}. Within the spectrum of PEFT strategies~\cite{houlsby2019parameter,aghajanyan2020intrinsic,hulora2022,edalati2022krona,wang2022adamix,gheini2021cross,zaken2022bitfit,guo2020parameter,sung2021training,ansell2022composable,lester2021power,li2021prefix,vu2022spot,he2021towards,mao2021unipelt,karimi2021compacter,liu2022few,sung2022lst,chen2022parameter,jia2022visual,chen2022adaptformer,zhang2022neural,jie2023fact,lian2022scaling,luo2023towards}, reparameterization-based variants~\cite{aghajanyan2020intrinsic,hulora2022,edalati2022krona} stand out as most closely related to our approach. For instance, LoRA~\cite{hulora2022} adapts a pretrained weight matrix by incorporating a low-rank perturbation, specifically the product of two learnable low-rank matrices, thus delivering strong results in NLP applications. While LoRA adopts a fixed-rank scheme across layers, subsequent work~\cite{zhang2022adaptive,valipour2022dylora,zhang2023increlora} has explored dynamic rank assignment tailored to different layers for more efficient parameter use. The ease of implementing low-rank reparameterizations has made such methods widespread~\cite{dettmers2023qlora,zi2023delta,chavan2023one}. Drawing inspiration from studies of hyperspherical energy and its link to model generalization~\cite{liu2018learning,liu2021orthogonal}, \cite{qiu2023controlling} introduced orthogonal finetuning (OFT), which learns an orthogonal transformation within each layer, thereby offering both improved generalization and consistently stable training outcomes relative to LoRA. However, OFT often incurs a higher count of trainable parameters than LoRA, motivating further advancements in its parameter efficiency. Furthermore, OFT's effectiveness beyond text-to-image diffusion model adaptation~\cite{qiu2023controlling} has not yet been thoroughly explored. BOFT advances this area by leveraging butterfly factorization to reduce parameter count, which in turn allows us to validate orthogonal finetuning across diverse adaptation tasks.

\textbf{Butterfly structures}. The classic radix-2 Cooley-Tukey algorithm decomposes the $N$-point discrete Fourier transform recursively into two subproblems of size $\frac{N}{2}$, resulting in a structural pattern known as the butterfly, which can be represented as the product of a sequence of sparse matrices (referred to collectively as a butterfly matrix). These butterfly matrices have previously been deployed to represent orthogonal matrices to sidestep pivoting in Gaussian elimination and to boost computational efficiency~\cite{parker1995random}, to facilitate stable training in recurrent neural networks~\cite{jing2017tunable}, and for kernel matrix approximation~\cite{munkhoeva2018quadrature}. Additionally, fast linear transformations based on butterfly parameterizations have been developed in~\cite{dao2019learning,dao2019kaleidoscope}, while~\cite{chen2022pixelated,dao2022monarch} integrate butterfly matrices into neural network architectures for more resource-efficient training. Their utility extends to accelerating matrix-vector multiplication~\cite{michielssen1996multilevel,o2010algorithm}, approximating large, structured matrices~\cite{li2015butterfly}, and optimizing network data transmission~\cite{klasing1994broadcasting,li2003linear,rajasingh2016transmission}. Distinct from previous studies, our work centers on exploiting butterfly structures to substantially improve OFT’s parameter efficiency.

\section{An Information Transmission View on Orthogonal Finetuning}

We begin by outlining the foundations of OFT. Rather than directly modifying the pretrained weight matrix, OFT introduces a reparameterization where the new weight matrix is represented as the product of a frozen pretrained weight matrix and a trainable orthogonal matrix. This approach contrasts with LoRA, which implements updates through the addition of a low-rank matrix, whereas OFT applies modifications by multiplying with an orthogonal matrix. For enhanced parameter efficiency, LoRA employs a low-rank configuration, while OFT initially utilizes a block-diagonal orthogonal form~\cite{qiu2023controlling}; efficiency is improved as the diagonal blocks decrease in size. Figure~\ref{fig:comp_lora} provides a visual juxtaposition of these strategies. The underlying rationale for utilizing an orthogonal transformation to adapt the weight matrix lies in maintaining the pair-wise angular relationships among neurons~\cite{liu2018learning,liu2021orthogonal,qiu2023controlling}, thereby allowing the pretrained semantic representations to be retained. Specifically, for a pretrained linear transformation $\bm{W}^0\in\mathbb{R}^{d\times n}$, OFT learns an orthogonal matrix $\bm{R}\in\mathbb{R}^{d\times d}$, altering the layer’s forward computation from $\bm{z}=(\bm{W}^0)^\top\bm{x}$ to $\bm{z}=(\bm{R}\bm{W}^0)^\top\bm{x}$, where $\bm{x}\in\mathbb{R}^d$ denotes the input and $\bm{z}\in\mathbb{R}^n$ is the output. OFT initializes $\bm{R}$ as the identity matrix to achieve zero initialization. To retain the orthogonality constraint during finetuning, we use the Cayley parameterization following~\cite{qiu2023controlling,liu2021orthogonal}: $\bm{R}=(\bm{I}+\bm{Q})(\bm{I}-\bm{Q})^{-1}$, where $\bm{Q}$ is skew-symmetric, i.e., $\bm{Q}=-\bm{Q}^\top$. For further parameter savings, $\bm{R}$ adopts a block-diagonal construction, parameterized as $\text{diag}(\bm{R}_1,\bm{R}_2,\cdots,\bm{R}_r)$, where each $\bm{R}_i\in\mathbb{R}^{b\times b}$ is a compact orthogonal block and $br=d$. Although this block-diagonal design enables parameter-efficiency, it imposes a structural assumption: the neuron dimensions (that is, each column in $\bm{W}^0$) are partitioned into $r$ subsets, and each subset is separately transformed by a different orthogonal matrix. While this structure works well in practice, grouping neuron dimensions merely by their indices lacks conceptual justification, highlighting the appeal of dense orthogonal matrices. This observation leads to a central question: \emph{Is it possible to realize a dense orthogonal matrix while retaining the benefits of parameter-efficiency?}

To explore this issue, we suggest constructing a dense orthogonal matrix as a product of several sparse orthogonal matrices. For a consistent yet accessible understanding of the sparsity structures in orthogonal matrix decompositions, we reinterpret the task of generating a dense orthogonal matrix in OFT through the lens of information transmission. Specifically, forming a dense matrix $\bm{R}\in\mathbb{R}^{d\times d}$ as the product $\bm{R}=\bm{B}_m\bm{B}_{m-1}\cdots\bm{B}_1$, where each $\bm{B}_i$ is square, can be conceptualized as the process of information propagation across a $d\times (m+1)$ grid, as depicted in Figure~\ref{fig:info_trans}. This perspective is motivated by noting that a dense $d$-by-$d$ matrix signifies full connectivity from an initial set of $d$ nodes to another. The interpretation is straightforward: for matrix $\bm{R}$, if $R_{ij}=0$, the $j$-th input cannot influence the $i$-th output, whereas a nonzero $R_{ij}$ enables information flow. Thus, representing $\bm{R}$ by sequential multiplication of matrices $\bm{B}_m$ through $\bm{B}_1$ can be seen as successive information exchanges governed by the structure of each $\bm{B}_i$. The information moves layer by layer: starting with $\bm{B}_1$ and culminating in $\bm{B}_n$. An example is shown in Figure~\ref{fig:info_trans} with the factorization $\bm{R}=\bm{B}_5\bm{B}_4\bm{B}_3\bm{B}_2\bm{B}_1$, displaying their sparsity patterns and associated graphs. The graph results from unfolding the multiplication sequence: each matrix $\bm{B}_i$ serves as a map connecting nodes between two consecutive layers ($i$-th and $(i+1)$-th). Specifically, element $(j_1,j_2)$ in $\bm{B}_i$ signifies a directed connection from node $j_2$ in layer $i$ to node $j_1$ in layer $i+1$, where zeros indicate the absence of a link. For $\bm{B}_5\bm{B}_4\bm{B}_3\bm{B}_2\bm{B}_1$ to yield a dense outcome, each node at the starting layer must have a path to every node in the sixth layer. If we only examine $\bm{R}=\bm{B}_4\bm{B}_3\bm{B}_2\bm{B}_1$, corresponding to pathways from the first to fifth layer, it becomes evident that information from node $1$ fails to reach node $3$, which manifests as a zero entry at $(3,1)$ in the product’s sparsity pattern. The same correspondence appears when analyzing the factorization chains $\bm{R}=\bm{B}_3\bm{B}_2\bm{B}_1$ and $\bm{R}=\bm{B}_2\bm{B}_1$. Broadly, in order for a matrix $\bm{R}\in\mathbb{R}^{d\times d}$ to be fully dense, the constituent matrices $\bm{B}_m,\ldots,\bm{B}_1$ must collectively realize a directed edge set across a $d\times(m+1)$ grid, where edges only connect nodes at adjacent layers (i.e., adjacent columns), guaranteeing that all starting nodes have transmission paths to every endpoint node at the final level.

The information transmission framework provides valuable insight into understanding the sparsity patterns of factorization matrices in OFT. As shown in Figure~\ref{fig:blk_diag}, the original OFT produces a block-diagonal configuration for $\bm{R}$. While this design reduces the total number of trainable parameters, it falls short in producing a densely connected matrix $\bm{R}$. Our objective is to assemble a dense orthogonal matrix from $m$ sparse orthogonal matrices, all while minimizing the number of effective trainable parameters. From the perspective of information transmission, two main criteria guide this process: \emph{(i)} \textbf{dense connectivity}, which requires that each node at the first layer connects via at least one path to every node at the final layer, and \emph{(ii)} \textbf{minimum free edges}, stipulating that—within the orthogonality constraints—the sum of all edges should be as limited as possible. The requirement of orthogonality imposes nuanced constraints on edges linking successive levels. Specifically, to ensure each matrix $\bm{B}_i$ is full-rank—a prerequisite for orthogonality—a bijective mapping involving $d$ edges between nodes at level $i$ and level $(i=1)$ must be maintained, meaning there are at least $d$ edges connecting adjacent layers (see Figure~\ref{fig:info_trans} for an example with 4). These $d$ connections are mandated by orthogonality and are exempt from edge counts, as their corresponding entries are not trainable (for instance, in a $d\times d$ orthogonal matrix with $d$ non-zero elements, these entries must be $\pm 1$). Because orthogonality reduces the parameter count compared to unrestricted matrices, it also induces further dependencies among edges. Take, for example, a $2\times 2$ orthogonal matrix: a zero at the $(1,1)$ entry necessarily enforces a zero at the $(2,2)$ entry (\emph{i.e.}, the absence of one edge may entail the absence of another). Consequently, for any valid configuration of edge connections, orthogonality may sometimes necessitate adding or eliminating specific edges. The information transmission approach, which centers on visualizing the distribution of non-zero elements in the resulting orthogonal matrix, proves particularly advantageous for OFT, since our primary focus is on the non-zero, trainable entries of $\bm{R}$; their actual values are inconsequential.

A straightforward approach to densely connecting two hierarchical levels requires $\mathcal{O}(d^2)$ edges (that is, a dense orthogonal matrix), corresponding to $d^2-d$ edges in total (for instance, when $d=4$, there are 12 edges involved). Figure~\ref{fig:info_trans} illustrates an example of a possible matrix factorization, which utilizes $10$ edges overall—fewer than are necessary for a dense orthogonal connection. This analytical framework provides a graphical means to assess the parameter-efficiency of OFT, facilitating the development of suitable factorizations. Drawing on the structure known as the butterfly graph from the Cooley-Tukey algorithm~\cite{cooley1965algorithm}, we note that such a topology permits an efficient dense mapping between $d$ source and $d$ receiver nodes with only $\mathcal{O}(d\log d)$ edges. To illustrate, while the structure in Figure~\ref{fig:info_trans} achieves dense connectivity with $10$ edges, the butterfly architecture accomplishes this with just $8$. The following section details how incorporating the butterfly structure can enhance parameter efficiency.

\section{Orthogonal Parameterization by Butterfly Factorization}

The butterfly structure, initially employed in the Cooley-Tukey algorithm, facilitates efficient computation of the fast Fourier transform. In the context of the Fourier transform, localized alterations within the frequency space result in widespread effects across the spatial domain, paralleling our objective: ensuring that each node at the initial layer can relay information to every node at the terminal layer. The butterfly configuration has also become prevalent as a networking topology~\cite{solihin2015fundamentals,li2011linear}, optimized for streamlined communication. Let $k\geq 2$ be a power of $2$. The \emph{butterfly factor} $\bm{B}^F(k)$ is then given by:

\begin{equation}\label{eq:but_factor}
\begin{aligned}
    \bm{B}^F(k)=\begin{bmatrix}
    \text{diag}(\bm{d}_1) & \text{diag}(\bm{d}_2) \\
    \text{diag}(\bm{d}_3) & \text{diag}(\bm{d}_4)
    \end{bmatrix}\in\mathbb{R}^{k\times k},
\end{aligned}
\end{equation}

where each $\bm{d}_i\in\mathbb{R}^{\frac{k}{2}}$ for $i=1,\ldots,4$ denotes a vector. Taking $d=2^N$, the $d$-dimensional \emph{butterfly matrix} $\bm{B}(d)\in\mathbb{R}^{d\times d}$ is then defined recursively as:

\begin{equation}
\begin{aligned}
    \!\!\bm{B}(d)=\tilde{\bm{B}}(d,d)\!\cdot\!\begin{bmatrix}
    \bm{B}_1(\frac{d}{2})\!\!\!\!\! & \bm{0} \\
    \bm{0}\!\!\!\!\! &\bm{B}_2(\frac{d}{2})
    \end{bmatrix}= \tilde{\bm{B}}(d,d)\tilde{\bm{B}}(d,\frac{d}{2})\cdots\tilde{\bm{B}}(d,2),
\end{aligned}
\end{equation}

Here, $\bm{B}_1(\frac{d}{2})$ and $\bm{B}_2(\frac{d}{2})$ represent two butterfly matrices, each of dimension $\frac{d}{2}$. Next, we introduce the concept of a \emph{butterfly component}, denoted as $\thickmuskip=2mu \medmuskip=2mu \tilde{\bm{B}}(d,k)=\text{diag}(\bm{B}^F_1(k),\cdots,\bm{B}^F_{d/k}(k))$, which is a $d\times d$ block-diagonal matrix composed of blocks of size $k$. Each block $\bm{B}^F_i(k), \forall i$ refers to a butterfly factor as defined in Equation~\ref{eq:but_factor}. Equipped with this, we proceed to leverage the butterfly matrix as a means to parameterize an orthogonal matrix. This can be accomplished by ensuring that each multiplicative component $\tilde{\bm{B}}(d,k), \forall k$ that makes up the butterfly matrix $\bm{B}(d)$ is orthogonal. 

We begin by examining the particular case of the block-diagonal matrix $\tilde{\bm{B}}(d,2)$, which has blocks of size $2$. Orthogonality of $\tilde{\bm{B}}(d,2)$ can be straightforwardly secured by parameterizing each $2 \times 2$ block through the Cayley transform (equivalently, via $2$-dimensional rotation), following the methodology outlined in \cite{liu2021orthogonal,qiu2023controlling}. Since the sparsity structure, i.e., the non-zero entries, of every butterfly component can be obtained by permuting the non-zero structure of $\tilde{\bm{B}}(d,2)$, all butterfly components can, as a result, be conveniently expressed as orthogonal matrices. This approach enables an efficient orthogonal matrix parameterization constructed from numerous $2\times 2$ orthogonal matrices.

Expanding upon this idea, inspired by \cite{chen2022pixelated}, we define a generalized version, called the block butterfly component $\tilde{\bm{B}}^b(d,k)$, wherein every scalar element in $\bm{d}_i, \forall i$ is replaced with a $b\times b$ matrix. To make the block butterfly component $\tilde{\bm{B}}^b(d,2)$ orthogonal, we enforce each $2b\times 2b$ block to be orthogonal. The block-level permutation of non-zero patterns for other butterfly components $\tilde{\bm{B}}^b(d,k), k>2$ mirrors the structure of $\tilde{\bm{B}}^b(d,2)$, making it possible to analogously enforce orthogonality for these cases. Collectively, these constructions form the basis of the BOFT forward pass, which is given by

\begin{equation}
    \begin{aligned}\nonumber
    \bm{z}=\big(\bm{R}(m,b)\cdot\bm{W}^0\big)^\top\bm{x},~~~~\text{s.t.}~\bigg{\{}\bm{R}(m,b)=\prod_{i=1}^m \tilde{\bm{B}}^b_{(d,i)}~~\&~~\underbrace{\big(\tilde{\bm{B}}^b_{(d,j)}\big)^\top\tilde{\bm{B}}^b_{(d,j)}=\tilde{\bm{B}}^b_{(d,j)}\big(\tilde{\bm{B}}^b_{(d,j)}\big)^\top\!=\bm{I}_d}_{\forall j\in [1, m]}\bigg{\}},
    \end{aligned}
\end{equation}

For brevity, we rewrite $\tilde{\bm{B}}^b(d,2^{m - i+1})$ as $\tilde{\bm{B}}^b_{(d,i)}$, and $\bm{I}_d$ denotes the $d$-dimensional identity matrix. The matrix $\bm{R}(m,b)\in\mathbb{R}^{d\times d}$ is an orthogonal matrix that results from the product of several orthogonal butterfly blocks. We refer to BOFT using $\bm{R}(m,\frac{b}{2})$ as BOFT($m$,$b$), assuming $b\geq 2$. When $\thickmuskip=2mu \medmuskip=2mu m=1$, BOFT($1$,$b$) simplifies to the block-diagonal OFT~\cite{qiu2023controlling} with block size $b$. Furthermore, BOFT($1$,$d$) coincides with the original OFT~\cite{qiu2023controlling} utilizing an unconstrained, full-sized orthogonal matrix. When employing $\log \frac{2d}{b}$ layers and block size $b$, that is, BOFT($\log \frac{2d}{b}$,$b$), the resulting matrix $\bm{R}$ is a block butterfly matrix $\bm{B}^{\frac{b}{2}}(d)$, providing a dense orthogonal transformation. More generally, BOFT($m$,$b$) incorporates $\frac{1}{2}(b-1)dm$ effective parameters to adapt a linear layer with dimensions $d\times n$. Specializing to butterfly matrices ($m = \log d$, $b = 2$), BOFT's parameter count scales as $\mathcal{O}(d\log d)$. Comparatively, the original OFT relying on a dense orthogonal matrix has $\mathcal{O}(d^2)$ parameters, while the block-diagonal OFT with $r$ blocks has parameter complexity $\mathcal{O}(bd)$. Notably, the standard OFT requires $b=d$ to achieve a dense orthogonal transformation, whereas BOFT can do so for arbitrary $b$.

\textbf{BOFT Identity Initialization}. Fine-tuning strategies often retain the parameters from the pretrained network to minimize divergence during adaptation. As an example, LoRA utilizes a zero initialization for its additional low-rank matrices. For BOFT, initialization sets every butterfly component to the identity matrix (that is, setting the skew-symmetric matrix in the Cayley parameterization to zero).

\textbf{BOFT Multiplicative Dropout}. LoRA~\cite{hulora2022} incorporates a Dropout mechanism within its low-rank adaptations to mitigate overfitting risks. While classic Dropout~\cite{srivastava2014dropout} aligns naturally with LoRA, it is incompatible with BOFT due to the multiplicative nature of BOFT's parameter updates. Thus, we introduce a multiplicative Dropout scheme designed for BOFT: the orthogonal matrix $\bm{R}(m,b)$, constructed from $m$ butterfly factors, can be rearranged into a $2b\times 2b$ block-diagonal structure. The proposed multiplicative Dropout randomly chooses a fraction $p_1$ of butterfly components and a proportion $p_2$ of the diagonal blocks within each component, substituting those selected blocks with identity matrices.

\section{Intriguing Insights and Discussions}

\textbf{Expressivity of BOFT}. While the butterfly structure, combined with appropriate permutations, can exactly replicate several classical fast linear transforms such as the fast Fourier transform and the Hadamard transform~\cite{dao2019learning,dao2019kaleidoscope}, it is still an open question how closely our orthogonal butterfly matrix can approximate an arbitrary orthogonal matrix. To investigate this, we simulate the process of approximating a randomly generated dense orthogonal matrix of dimension $512\times 512$~\cite{anderson1987generation}. The findings, depicted in Figure~\ref{fig:para_eff}, are based on averaging across 10 different random seeds. On these plots, the vertical axis indicates the error of the approximation, while the horizontal axis shows the count of effective, trainable parameters. Curves sharing the same color correspond to BOFT with identical block sizes; the curve's leftmost point represents the error associated with BOFT($1$,$b$), that is, the standard block-diagonal OFT of block size $b$. In general, BOFT attains superior parameter efficiency compared to OFT. For instance, BOFT($9$,$2$) is more expressive than BOFT($1$,$16$) while utilizing far fewer parameters. Moreover, BOFT configurations with smaller $b$ and larger $m$ tend to be more efficient; e.g., BOFT($6$,$4$) reaches a similar error as BOFT($2$,$16$) but with notably fewer parameters. Collectively, these results suggest that the butterfly matrix encapsulates a more structured subset within the orthogonal group, in contrast to the block-diagonal structure, thereby establishing that BOFT possesses provably greater expressivity than OFT with equivalent block size.

\begin{theorem}[Expressivity of BOFT]\label{thm:exp_boft}
    For the same block size, BOFT demonstrates strictly greater expressivity compared to OFT. To approximate all orthogonal matrices of dimension $d$, products of butterfly matrices of the form $\bm{B}_{d-1,1}(d)\bm{B}^\top_{d-1,2}(d)\cdots\bm{B}_{1,1}(d)\bm{B}^\top_{1,2}(d)$—with each $\bm{B}_{i,j}(d)$ being a butterfly matrix for every $i$ and $j$—can be employed.
\end{theorem}

Theorem~\ref{thm:exp_boft} provides a straightforward extension for BOFT: the final orthogonal matrix can be generalized as $\thickmuskip=2mu \medmuskip=2mu \bm{R}^G(m_1,b_1,m_2,b_2,l)=\bm{R}_{l,1}(m_1,b_1)\bm{R}_{l,2}^T(m_2,b_2)\cdots\bm{R}_{1,1}(m_1,b_1)\bm{R}_{1,2}^T(m_2,b_2)$, where $\bm{R}_{i,j}^T(m,b)$ are orthogonal matrices employed in the BOFT construction. Specifically, by choosing $m_1=m_2=\log d$, $b_1=b_2=2$, and $l=d-1$, this matrix, $\bm{R}^G(m_1,b_1,m_2,b_2,l)$, is capable of spanning the entire orthogonal group. This composition aligns with the kaleidoscope hierarchy~\cite{dao2019kaleidoscope}. Nevertheless, increased expressivity does not universally guarantee superior performance during finetuning: for example, although full finetuning is universally expressive, it often results in suboptimal outcomes. Thus, balancing expressivity with regularization emerges as central for achieving effective model generalization. The structural priors incorporated by BOFT expand the space of tunable parameters, aiding in discovering a more advantageous compromise between expressivity and regularization.

\textbf{Spectral properties}. Compared to LoRA, orthogonal finetuning more effectively maintains desirable spectral characteristics, as it exactly retains the spectral norm of the original pretrained weight matrix $\bm{W}^0$. To illustrate, consider the singular value decomposition $\bm{W}^0 = \bm{U}\bm{\Sigma}\bm{V}^\top$, where $\bm{U}$ and $\bm{V}$ denote orthogonal matrices and $\bm{\Sigma}$ is a diagonal matrix containing singular values. When either OFT or BOFT is applied, an orthogonal matrix $\bm{R}$ multiplies on the left, yielding updated weights of the form $\bm{R}\bm{U}\bm{\Sigma}\bm{V}^\top$. This operation leaves the maximal singular value, or equivalently, the spectral norm of $\bm{W}^0$, unchanged. Prior work demonstrates that maintaining the spectral norm in this manner facilitates more stable optimization and improves generalization performance~\cite{miyato2018spectral,yoshida2017spectral}. Additional theoretical insights regarding these properties are available in Appendix~\ref{app:sn_property}.

\textbf{Orthogonal finetuning as learning bilinear similarity}. In BOFT, the process can be recast as learning a bilinear similarity metric of the form $\bm{w}^0_i\bm{R}\bm{x}$, with $\bm{w}^0_i$ denoting the $i$-th column (representing a neuron) of $\bm{W}^0$. Here, the bilinear matrix $\bm{R}$ is constrained to be orthogonal, imposing significant structure. This formulation naturally relates to classical problems such as distance metric learning~\cite{xing2002distance} and analysis of bilinear forms~\cite{roman2005advanced}.

\textbf{Inductive bias and generalization in BOFT}. The matrix $\bm{R}(m,b)$ used in BOFT often parameterizes a specific, structured subset of the orthogonal group, thereby restricting the space of functions BOFT can represent and inducing a distinct inductive bias. The structure inherent in butterfly factorization, shared with many well-known linear operators (such as the discrete Fourier, sine/cosine, and Hadamard transforms~\cite{dao2019kaleidoscope}), is argued to encourage better generalization. Furthermore, the approach’s reliance on sparse matrix factorizations introduces additional forms of implicit inductive bias, as previously discussed in related literature~\cite{gunasekar2017implicit,li2018algorithmic,liu2019neural}.

\textbf{Contrast with butterfly-based sparse training}. Several prior works~\cite{dao2019learning,dao2019kaleidoscope,chen2022pixelated,dao2022monarch} have investigated sparse training techniques that utilize the butterfly parameterization. These approaches generally emphasize reparameterizing neural network weight matrices through butterfly structures and subsequently training the models from initialization. In particular, \cite{dao2022monarch} explores a finetuning paradigm in which pretrained weights are initially mapped onto a form of butterfly matrices, after which the projected parts are finetuned for target tasks. In contrast, BOFT adopts a fundamentally different finetuning methodology, leveraging layer-shared weight matrices to transform the weights.

\input{rebuttal-results}

\section{Concluding Remarks and Limitations}

In this study, we introduce Orthogonal Butterfly, a broadly applicable approach for parameter-efficient finetuning of foundation models that leverages the butterfly structure. Central to our method is the observation that superior parameter efficiency is attainable by representing a dense orthogonal matrix as a product of several sparse orthogonal matrices. To facilitate the discovery of appropriate matrix decompositions, we put forward a graphical information transmission perspective. Through this lens, we demonstrate that the butterfly structure is especially well-suited for our objective of sparse orthogonal matrix factorization. Our empirical evaluation highlights the effectiveness of BOFT when finetuning a range of models, including large language models, large-scale vision architectures, and text-to-image generators. Furthermore, experimental results indicate that BOFT outperforms other general-purpose finetuning strategies.

Nonetheless, BOFT is not without drawbacks. Owing to the structure in which the resulting orthogonal matrix is composed via a product of multiple orthogonal matrices, training incurs somewhat higher runtime overhead relative to OFT. Addressing the training efficiency of BOFT is an open area for future research. On the positive side, after the finetuning phase, the orthogonal matrices learned by BOFT can be fused directly into the original pretrained model, incurring no additional latency during inference. Another open question is whether the butterfly network truly represents the most effective structure for information transmission. Our framework for modeling information transmission suggests potential avenues for borrowing insights from fields such as computer networking, where network topologies are extensively analyzed for their information dissemination properties. We believe that this opens the door for the exploration of alternative network structures that could yield more efficient compositions for dense orthogonal matrices.

\end{document}